\documentclass[11pt]{article}

% Standard packages for arXiv compatibility
\usepackage[utf8]{inputenc}
\usepackage{amsmath,amssymb}
\usepackage{graphicx}
\usepackage{booktabs}
\usepackage{hyperref}
\usepackage{natbib}
\usepackage{xcolor}
\usepackage{microtype}
\usepackage{placeins}  % provides \FloatBarrier

% Page setup
\usepackage[margin=1in]{geometry}

% Math commands
\newcommand{\R}{\mathbb{R}}
\newcommand{\N}{\mathbb{N}}
\newcommand{\Z}{\mathbb{Z}}
\newcommand{\E}{\mathbb{E}}
\newcommand{\Prob}{\mathbb{P}}

% Theorem environments
\usepackage{amsthm}
\newtheorem{theorem}{Theorem}[section]
\newtheorem{lemma}[theorem]{Lemma}
\newtheorem{proposition}[theorem]{Proposition}
\newtheorem{corollary}[theorem]{Corollary}
\newtheorem{definition}[theorem]{Definition}
\newtheorem{remark}[theorem]{Remark}

\title{Your Paper Title Here}

\author{
  First Author\thanks{Equal contribution.} \\
  Department of Science\\
  University Name\\
  \texttt{first@university.edu} \\
  \And
  Second Author\footnotemark[1] \\
  Department of Engineering\\
  Another University\\
  \texttt{second@another.edu}
}

\date{}

\begin{document}

\maketitle

\begin{abstract}
Your abstract goes here. It should be a single paragraph summarizing the
problem, approach, key results, and significance of the work. Keep it
under 250 words for most venues.
\end{abstract}

\section{Introduction}
\label{sec:intro}

Introduce the problem and motivation here. Cite relevant prior work
\citep{example2024} and clearly state the contributions of this paper.

\paragraph{Contributions.} We summarize our main contributions:
\begin{enumerate}
    \item First contribution.
    \item Second contribution.
    \item Third contribution.
\end{enumerate}

\section{Related Work}
\label{sec:related}

Discuss relevant prior work here, organized by topic.

\section{Method}
\label{sec:method}

Describe your approach here.

\subsection{Problem Setup}
\label{sec:setup}

Formally define the problem.

\subsection{Proposed Approach}
\label{sec:approach}

Detail the proposed method or algorithm.

\section{Experiments}
\label{sec:experiments}

\subsection{Experimental Setup}

Describe datasets, baselines, metrics, and implementation details.

\subsection{Results}

Present your main results.

% Example table
% \begin{table}[t]
% \centering
% \caption{Main results comparison.}
% \label{tab:results}
% \begin{tabular}{lcc}
% \toprule
% Method & Metric 1 & Metric 2 \\
% \midrule
% Baseline 1 & 0.00 & 0.00 \\
% Baseline 2 & 0.00 & 0.00 \\
% \textbf{Ours} & \textbf{0.00} & \textbf{0.00} \\
% \bottomrule
% \end{tabular}
% \end{table}

\subsection{Analysis}

Provide ablation studies and analysis of results.

\section{Discussion}
\label{sec:discussion}

Discuss limitations and broader impact.

\section{Conclusion}
\label{sec:conclusion}

Summarize the paper and suggest future directions.

\section*{Acknowledgments}

Acknowledge funding sources and helpful discussions.

\bibliographystyle{plainnat}
\bibliography{references}

% Appendix (optional)
% \appendix
% \section{Additional Experiments}
% \label{app:experiments}

\end{document}
